\documentclass[../main.tex]{subfiles}
\begin{document}
%==========================================TIÊU ĐỀ TẬP========================================
\begin{table}[h]
  \begin{adjustwidth}{-1cm}{}
    \begin{tabular}{>{\centering\arraybackslash}p{6cm}|>{\raggedright\arraybackslash}p{12.5cm}}
      \multirow{5}{6cm}
      % Cột bên trái: ÉPISODE và số tập
      {\\[-40pt]
        \flaregothic\fontsize{20pt}{20pt}\selectfont \centering
        \color{eptype}ÉPISODE\color{black}\\
        \burbank\fontsize{72pt}{72pt}\selectfont
        \color{epnum}55\color{black}\\[6pt]
        \cafeteria\fontsize{20pt}{20pt}\selectfont\color{epnum}(4-13)\color{black}
        \\[18pt]
        \flaregothic\fontsize{18pt}{18pt}\selectfont
        \color{epattr}ÉPISODE PERDU\color{black}
      }
       &
      % Dòng thứ nhất: tên tập tiếng Nhật
      {\fotskip\fontsize{18pt}{18pt}\selectfont \ruby{最}{さい}\ruby{強}{きょう}の\ruby{挨}{あい}\ruby{拶}{さつ}を\ruby{決}{き}meろ\ruby{☆}{ほし}\ruby{挨}{あい}\ruby{拶}{さつ}\ruby{王}{おう}\ruby{選}{せん}\ruby{手}{しゅ}\ruby{権}{け}！
      } \\[6pt]
      % Dòng thứ hai: tên tập tiếng Anh
       & {\cafeteria\fontsize{18pt}{18pt}\selectfont A Carnival of Greetings!}                                                                                                                      \\[4pt]
      % Dòng thứ ba: tên tập tiếng Việt
       & {\vnmsans\fontsize{18pt}{18pt}\selectfont Giải vô địch Vua chào hỏi}                                                                                                                       \\[8pt]
      % Dòng thứ tư: Nhân vật xuất hiện
       & {\caviar{\fontsize{14pt}{14pt}\selectfont Nhân vật: \emp{kuon1}, \emp{achan}, \emp{matsuri}, \emp{ayame}, \emp{rushia}}}
      \\[8pt]
      % Dòng cuối cùng: Thời gian diễn ra của tập (thường sẽ là ngày phát hành)
       & {\continuum\fontsize{16pt}{16pt}\selectfont Ngày phát hành: 24/05/2020}                                                                                                                    \\[18pt]
    \end{tabular}
  \end{adjustwidth}
\end{table}
%=========================================TRUYỆN CHÍNH========================================
\color{black}

\txtbx[2]{{\color{vol} \uline{Tóm tắt:}}} Một đêm muộn nọ, Matsuri khiến cả văn phòng đứng hình khi hồ hởi cất tiếng chào ``Ohayo'' (Chào buổi sáng) dù ngoài trời đã tối mịt. Sự mâu thuẫn đầy thú vị này đã khơi mào cho một cuộc thi sáng tạo câu chào độc lạ vô tiền khoáng hậu giữa Matsuri, Ayame và Rushia. Từ màn chào mang đậm sắc màu tôn giáo, phép thuật ``bắt chước'' câu chào của các thành viên khác cho đến chiêu thức thôi miên ``Hypno'' của Ayame, cuộc thi kết thúc bằng một ván thôi miên tập thể dở khóc dở cười đưa cả nhóm chìm sâu vào giấc ngủ ngon lành cho đến sáng hôm sau.

\begin{center}
  \uline{(Người viết tập này có sử dụng bản dịch từ \href{[https://www.youtube.com/@TomangclipHololivevedich](https://www.youtube.com/@TomangclipHololivevedich)}{Tớ mang clip \hll\ về dịch - TMCHVD}.)}
\end{center}

Matsuri rạng rỡ đẩy mở toang cánh cửa phòng họp với những bước chân tràn đầy sự hứng khởi, chất giọng vui tươi, lanh lảnh vang lên ngập tràn năng lượng tựa như tiếng chuông ngân vang: ``Chào buổi sáng cả nhà nhé!''

Câu chào hỏi vô cùng lệch pha đó của cô nàng lập tức khiến cho toàn bộ mọi người có mặt bên trong căn phòng khựng lại mất một nhịp. Ayame đang chăm chú lướt điện thoại ngơ ngác ngẩng đầu lên nhìn, Rushia chậm rãi quay người lại từ góc bàn họp, còn Kuon cũng đành phải tạm thời đặt bút dừng việc viết bản báo cáo tài chính dở dang. Ba ánh mắt ngập tràn sự hoang mang cùng hàng tá dấu chấm hỏi to đùng đồng loạt đổ dồn về phía cô nàng Cổ động.

Nàng Quỷ sừng và nàng Pháp sư khẽ rụt rè cất tiếng đáp lại lời chào của cô một cách ngập ngừng: ``Dạ, chào buổi sáng chị ạ\dots'' dẫu cho gương mặt của Rushia vẫn đang lộ rõ vẻ vô cùng lưỡng lự và bối rối.

Kuon thọc tay túi quần đáp lại một cách vô cùng máy móc theo thói quen phản xạ: ``Osu. Cơ mà khoan đã em ơi\dots'' Vừa dứt lời chào, ánh nhìn của cậu Tổng đã nhanh chóng chuyển dịch hướng thẳng ra ngoài khung cửa sổ kính.

Bầu trời ở phía bên ngoài lúc này đã được nhuốm một sắc màu xanh dương đậm đặc của màn đêm tĩnh mịch, vài ngọn đèn đường neon xung quanh khu nhà đã sớm được bật sáng rực rỡ từ bao giờ, thậm chí trong không gian vắng lặng thỉnh thoảng còn có thể nghe thấy rõ mồn một tiếng chim cú kêu rúc từng hồi ngoài rừng. Ayame chớp chớp mắt đầy hoang mang, dùng ngón tay chỉ thẳng về phía khung cửa sổ tối đen, ngập ngừng lên tiếng chất vấn: ``Nhưng mà Matsuri-senpai, ngoài trời bây giờ rõ ràng đã tối mịt tối mò thế kia rồi cơ mà? Chào buổi sáng cái tầm này nghe nó cứ có gì đó sai sai thế nào ấy.''

\img{4-13a}

Rushia cũng gật gù đồng tình lia lịa: ``Dạ đúng thế luôn đấy-nanodesu. Em thực sự cũng thấy cái câu chào của chị nó kỳ lạ và bất hợp lý lắm luôn á, Matsuri-senpai.''

Cô nàng Lễ hội Matsuri nghe thấy hai cô em phàn nàn liền chẳng hề bận tâm, cô nàng chỉ nở một nụ cười tươi rói tự mãn, thong thả lên tiếng giải thích để giải vây cho sự nghi hoặc của hai cô gái: ``Bọn em đúng là chẳng am hiểu sự đời gì cả! Ở đất nước Nhật Bản của tụi mình á, mọi người trong giới giải trí hay thần tượng thường có thói quen cất tiếng chào `Ohayo' (Chào buổi sáng) vào bất cứ lúc nào khi lần đầu tiên gặp mặt nhau trong ngày đấy nhé, bất kể là lúc đó trời đang sáng trưng hay đã tối mịt tối mò từ lâu rồi.''

Kuon nghe xong liền đặt chiếc bút bi xuống mặt bàn, khoanh hai tay trước ngực với vẻ mặt nghiêm túc lên tiếng chỉnh đốn lại lời ngụy biện của Matsuri: ``Đúng ra xét về mặt ngữ pháp thông thường thì cái tầm đêm muộn thế này phải cất tiếng chào 'Konbanwa' (Chào buổi tối) mới là chuẩn chỉ nhất em ạ. Nhưng mà cái văn hóa đặc thù của giới giải trí bên Nhật nó vốn dĩ đã có thói quen quy định như thế rồi, giờ hễ có đứa nào làm việc trong ngành mà đi ngược lại, chào kiểu khác thì tự dưng sẽ bị mọi người xung quanh coi là một kẻ kỳ cục, lập dị ngay ấy mà.''

Ayame gật gù cái rụp, như vừa sực nhớ ra một miền ký ức liên quan: ``À đúng rồi, yo! Em nghe anh nhắc mới nhớ ra đấy. Phía ban quản lý nhân sự ngày trước cũng từng căn dặn bảo với tôi cái quy tắc chào hỏi đặc thù này rồi. Hình như là đợt mới nhận việc, chị A-chan đã từng cất công giải thích cặn kẽ chuyện này cho em nghe một lần rồi thì phải.''

``Ừm, chính xác là như vậy đấy, chính chị ấy đã dặn dò thiết lập nội quy như thế mà,'' cậu con trai gật đầu xác nhận nhấn mạnh, đôi mắt vẫn nhìn Matsuri với vẻ nửa tin nửa ngờ.

\img{4-13b}

Rushia dẫu đã nghe giải thích xong xuôi nhưng gương mặt vẫn giữ nguyên sự dè dặt và ái ngại vô cùng, cô nàng khẽ lẩm bẩm trong miệng: ``Hóa ra là vậy á\dots\ Mà nghe cái thói quen chào hỏi này nó cứ ngược ngạo, lạ lùng làm sao ấy-nanodesu\dots''

Nhận thấy các cô em vẫn chưa thể hoàn toàn chấp nhận và thấu hiểu nổi cái nét văn hóa kỳ lạ này, cô nàng Lễ hội Matsuri liền chống nạnh ngang hông nở một nụ cười lớn đầy tự hào: ``Thế nên người ta mới bảo đó là cái nét văn hóa đặc thù độc đáo chứ sao nữa! Các em không cảm thấy cái sự mâu thuẫn đó nghe có vẻ vô cùng thú vị và kích thích hay sao chứ hả?''

Cậu Tổng quản lý nhướng một bên mày, gương mặt lộ rõ vẻ bất lực pha chút châm chọc mỉa mai: ``Thú vị thì công nhận nghe cũng có phần vui tai thật đấy\dots\ nhưng mà anh nghĩ là hoàn toàn chẳng có chuyện ai bước chân vào cái văn phòng này cũng có thể dễ dàng làm quen được với cái cách chào hỏi kiểu củ chuối thế này ngay từ đầu đâu. Đến cả một thằng rành rọt công nghệ như anh mà cái hồi đầu tiên đặt chân tới đây nhận việc nghe mấy đứa chào Ohayo lúc nửa đêm còn thấy vô lý đùng đoàng, bỡ ngỡ chết đi được nữa là.''

``Nếu đã như vậy thì tại sao ba chị em bọn mình không thử cùng nhau động não, sáng tạo ra một cái cách thức chào hỏi mới mẻ độc lạ mang đậm thương hiệu của văn phòng mình đi xem nào?'' nàng Quỷ sừng bỗng chợt nảy ra một ý tưởng vô cùng táo bạo và có phần hơi điên rồ.

Ngay lập tức, Matsuri liền gật đầu tán thành nhiệt liệt không chút do dự, cô nàng cười híp cả mắt đầy phấn khích: ``Ý kiến này nghe tuyệt vời ông mặt trời quá đi, Ojou! Chị cũng đang có chung một suy nghĩ sáng tạo giống hệt như em đấy! Nào, để xem rốt cuộc trong ba chúng ta, ai mới thực sự là người sở hữu cái cách chào hỏi độc đáo, sáng tạo và đỉnh cao nhất nhé! Bọn mình sẽ tiến hành tổ chức một cuộc thi quyết đấu và kết quả sẽ được phân định ngay lập tức thôi!''

Trong khi hai cô nàng đang vô cùng hừng hực khí thế rực lửa, thì cậu con trai Kuon cùng nàng Pháp sư Rushia đứng bên cạnh vẫn giữ nguyên cái thái độ có phần hơi e dè, ái ngại với cái ý tưởng điên rồ này. Nhưng rồi dưới sức ép của số đông, họ cũng đành phải cắn răng tặc lưỡi thử sức tham gia vào một cái sự kiện dường như còn mang tính chất kỳ quặc và hỗn loạn hơn cả cái phong cách chào hỏi Ohayo ban đầu của cô nàng Lễ hội kia gấp hàng trăm lần.

\ct

Quyết tâm giật giải vô địch, Matsuri liền nhanh nhảu chạy ra ngoài hành lang rồi đẩy cửa hăm hở bước vào phòng họp một lần nữa. Lần này, cô nàng dang rộng cả hai cánh tay ra hai bên, tạo dáng đứng vô cùng kiêu sa lộng lẫy hệt như một vị thiên sứ thánh thiện vừa mới giáng trần, khẽ nháy mắt cất chất giọng vô cùng ngọt ngào, trìu mến lên chào hỏi: ``Chào buổi chiều tốt lành nhé, hỡi con chiên non nớt, đáng yêu\footnote{Cách gọi đầy ẩn ý mô phỏng theo danh xưng ``con cừu non'' của những người theo đạo Công giáo. Tuy nhiên trong ngôn ngữ lóng, từ ngữ này còn mang một tầng ý nghĩa khác châm biếm, ám chỉ một kẻ ngây thơ, dễ bị dắt mũi và có phần khờ khạo.} của lòng chị!''

Kuon đưa tay lên lau vệt mồ hôi trên trán, gương mặt chẳng xệch cố gắng che giấu đi cái nét ngơ ngác cạn lời trước màn làm màu của cô nàng: ``Ờ\dots\ Nghe cái câu chào hỏi này của em nó mang đậm đặc sắc màu tôn giáo giáo phái hắc ám thế nào ấy em ạ\dots\ nghe rợn tóc gáy vãi.''

Ayame khẽ nhướng đôi lông mày tinh quái, dùng ánh mắt không mấy tin tưởng nhìn cô nàng vừa bước vào phòng kia, nhún vai nhận xét: ``Cơ mà chị ơi, em thấy tụi mình có cần thiết phải rườm rà thêm cái đoạn từ ngữ 'chào buổi chiều' vào làm gì cho mệt không thế hả chị?''

Rushia đứng bên cạnh cũng khẽ nghiêng nghiêng đầu, đôi mắt nheo lại tỏ vẻ vô cùng thắc mắc: ``Em cũng đồng ý với Ojou đấy ạ, em nghĩ là cái đoạn danh xưng 'con chiên non' nghe có vẻ hơi bị thừa thãi và sến sẩm quá mức rồi đấy-nanodesu\dots''

\img{4-13c1}

Cậu Tổng quản lý đứng một bên nghe hai cô em phân tích chỉ trích liền khẽ nhíu mày mỉa mai châm chọc: ``Nóói tóm lại là theo như cái lời bình luận sắc sảo của hai đứa tụi em thì thà cô nàng Matsuri đừng có mở mồm ra chào hỏi câu nào luôn ngay từ đầu thì nghe nó còn có vẻ ổn áp và đỡ tốn diện tích hơn đấy nhỉ.''

Nàng Lễ hội tức tối chống nạnh hai tay ngang hông, phồng má giận dỗi cãi lại đầy đắc ý: ``Hai cái đứa này đúng là chẳng có tí khiếu cảm thụ nghệ thuật nào cả! Thừa thãi ở cái chỗ nào cơ chứ hả? Các em đằng ấy không cảm thấy cái cách chào hỏi tôn giáo này nghe vô cùng thần thánh, trang trọng và mang đậm tính chuyên nghiệp đỉnh cao hay sao!?''

\ct

Quyết tâm gỡ gạc lại thể diện cho giới Oni, Ayame bất ngờ lấy đà chạy phóc vào bên trong phòng họp. Đầu tiên cô nàng cúi gằm mặt xuống đất tạo sự huyền bí, rồi bất thình lình ngẩng phắt gương mặt xinh xắn lên nhìn mọi người, hai cánh tay dang rộng ngang ra hai bên và dũng cảm đứng thăng bằng chỉ bằng độc một chiếc chân trái, tạo nên một tư thế tạo dáng trông vô cùng giống\dots\ một con vịt Donald đang tập đi thăng bằng. Cô nàng Oni tự tin cất chất giọng lảnh lót của mình lên dõng dạc chào hỏi: ``Mágico\dots\ Banana!\footnote{Nguyên văn: マジカル…バナナ！, tức là trò chơi nối chữ liên tưởng 'Quả chuối ma thuật' vô cùng nổi tiếng của Nhật Bản.}''

Matsuri đứng đối diện nghe thấy thế liền vô cùng nhạy bén bắt nhịp nhập cuộc ngay tắp lự. Cô nàng một tay giơ cao nắm đấm lên trời tạo dáng, tay còn lại thì ôm chặt lấy chú cừu bông mềm mại trong lòng, hào hứng đáp lời theo đúng quy luật trò chơi liên tưởng: ``Quả chuối ma thuật làm cho chị lập tức liên tưởng đến cái màu vàng rực rỡ!''

Rushia cũng hớn hở vỗ tay bôm bốp nhảy vào hùa theo cuộc chơi: ``Màu vàng rực rỡ lại làm cho em liên tưởng ngay đến quả chuối ngọt ngào nè-nanodesu!''

\img{4-13c2}

Kuon ngồi tựa người ở góc phòng họp chứng kiến cái màn nối chữ liên tưởng vô tri của ba cô nàng liền chắp hai tay lên trán đầy mệt mỏi, cất chất giọng nghiêm nghị lên tiếng chặn họng dứt khoát: ``Anh xin can mấy đứa nhé, nghe cái trò chơi nối chữ này nó hãm tài và dở hơi cám lợn thế nào ấy em ạ\dots\ Hoàn toàn chẳng có một chút xíu nào mang dáng dấp của một câu chào hỏi cả! Anh không duyệt cái phương án này đâu nhé! Mau mau đổi sang cái kiểu chào khác hộ anh cái đi!''

Nàng Oni nghe thấy lời phê bình phũ phàng liền bĩu môi tiếc nuối, phồng hai má lên phụng phịu trách móc: ``Mồ\~{}! Đúng là Kuon-senpai chẳng có một tí xíu nào am hiểu và thấu cảm được cái độ sáng tạo, nghệ thuật đỉnh cao của em cả mà! Anh làm thế là tự dưng làm em mất sạch sành sanh cả nguồn cảm hứng sáng tạo luôn rồi đấy nhé!''

\ct

Đến lượt của cô đồng Rushia ra sân trổ tài. Cô nàng nhẹ nhàng, uyển chuyển bước từng bước chân chậm rãi tiến vào bên trong phòng họp, gương mặt hướng thẳng lên bầu trời, dùng một bàn tay khéo léo che hờ đi nửa gương mặt của mình, đôi mắt nhắm nghiền lại đầy vẻ ma mị, huyền bí hệt như một pháp sư đang chuẩn bị làm lễ cầu hồn. Cô nàng cất chất giọng u ám, thì thầm đầy bí ẩn chào hỏi: ``Hỡi những linh hồn lạc lối ngoài kia ơi, có ai đang nôn nóng muốn được biết cái hình thù của Ma cây gớm ghiếc nó ngang dọc thế nào không nào\~{}?''

Matsuri đứng bên cạnh nghe xong câu thoại ma mị liền lập tức nhảy dựng người lên phấn khích. Cô nàng giơ một tay lên cao tạo hình như đang cầm một thanh bảo kiếm ánh sáng tưởng tượng trong tay, tay còn lại vẫn ôm khư khư chú cừu bông, dõng dạc hét lớn nhập vai: ``Tiêu diệt ác ma trừ gian diệt bạo! Tiêu diệt ác ma trừ gian diệt bạo thôi nào mọi người ơi!''

Nàng Quỷ sừng cũng quyết tâm không chịu thua kém người bạn, cô nàng liền phối hợp hành động bằng cách\dots\ dùng tay tung mạnh một quả chuối chín vàng lên không trung tạo hiệu ứng: ``Hồi sinh người chết cải tử hoàn sinh! Hồi sinh người chết cải tử hoàn sinh thôi nào!''

\img{4-13c3}

Kuon đưa tay lên vỗ trán thở dài thườn thượt một tiếng rõ lớn, cố gắng nén cơn buồn cười đang dâng trào trong bụng: ``Cái màn biểu diễn tâm linh gọi hồn kết hợp tung chuối trừ tà này thì anh xin phép giơ hai tay đầu hàng chịu chết luôn nha. Đóng kịch thế là quá đủ rồi, tiếp tục người tiếp theo lên sóng hộ anh cái đi!''

Nàng Cô đồng tóc xanh nghe thấy thế liền mở to hai mắt tròn xoe nhìn ba người, phồng má giận dỗi đầy ấm ức: ``Thật là bực mình quá đi mất, ai cho phép các người tự tiện chen ngang phá nát hết cái bầu không khí cảm xúc ma mị nghệ thuật của tôi thế hả!''

\ct

Bất thình lình, không một điềm báo trước, cánh cửa gỗ phòng họp bỗng nhiên bị đẩy bật mạnh ra. A-chan thình lình xuất hiện lộ diện ngay tại lối vào với một vẻ mặt vô cùng quyết tâm, dứt khoát và cất giọng lớn tiếng tuyên bố: ``Tôi đã dõng dạc bảo với mọi người là cái quy tắc chào hỏi Ohayo đó nó bắt buộc phải như vậy rồi mà!'' 

Hoặc nói một cách chính xác và khách quan hơn, cái lời tuyên bố quả quyết hùng hồn đó của cô nàng trợ lý thực chất là đang muốn ra lệnh yêu cầu mọi người hãy mau chóng dừng cái cuộc thi sáng tạo vô tri, ngớ ngẩn này lại để quay về làm việc cho xong deadline.

Cậu Tổng quản lý Kuon đang ngồi viết báo cáo nghe thấy tiếng động liền ngoảnh đầu lại nhìn, nhướng một bên mày mỉa mai nhìn người đồng nghiệp: ``Nhưng mà chị A-chan ơi, cái câu tuyên bố hùng hồn đó của chị hoàn toàn chẳng có một cái nét nào mang dáng dấp cấu trúc của một lời chào hỏi cả đâu nhé\dots\ Với cả, chị tự nhìn ra xung quanh đi, ở cái phòng này hình như đang chẳng có lấy một mống nào thèm để tâm lắng nghe lời chị phát biểu kìa.''

\begin{figure}[H]
  \centering
  \begin{subfigure}[t]{0.45\textwidth}
    \centering
    \includegraphics[width=8cm]{../../../../Image/Book4/4-13c4.png}
    \caption{A-chan lao vào với kiểu ``chào'' như đang thuyết phục mọi người dừng lại\dots}
  \end{subfigure}
  \quad
  \begin{subfigure}[t]{0.45\textwidth}
    \centering
    \includegraphics[width=8cm]{../../../../Image/Book4/4-13c5.png}
    \caption{\dots nhưng không ai thèm nghe cả.}
  \end{subfigure}
\end{figure}

Quả đúng thật sự là như vậy, ba cô nàng Matsuri, Ayame và Rushia lúc này dường như hoàn toàn chẳng thèm để tâm hay mảy may có chút phản ứng nào trước cái câu ``chào'' chặn họng không giống ai của A-chan cả. Cô nàng Rushia đang vô cùng bận rộn cố gắng rướn người giữ thăng bằng cho một chiếc đĩa phim đặt trên đỉnh đầu mình, miệng không ngừng khúc khích cất tiếng gọi dụ dỗ chú cừu bằng bông đang đặt dưới sàn nhà: ``Lại đây chơi với chị nào, chú cừu non-chan dễ thương của chị ơi\~{}''

Ayame thì khom lưng cúi thấp người xuống đất, tay cầm quả chuối chín vàng giơ giơ lên trước mõm cừu bông dụ dỗ: ``Ngoan nào cừu ơi, chị có quả chuối ngọt ngào thơm phức dành riêng cho cưng nè\~{}''

Matsuri đứng bên cạnh khoanh hai tay trước ngực với một vẻ mặt vô cùng tự hào và tự mãn, dõng dạc tuyên bố khẳng định chủ quyền: ``Mấy đứa cứ nằm mơ đi nhé, tuyệt đối chẳng bao giờ có chuyện mấy đứa có thể dùng mấy cái món đồ ăn tầm thường kia để dụ dỗ được chú cừu con yêu quý của lòng chị đi theo đâu nhé.''

Đứng ngoài quan sát cái cảnh tượng dở khóc dở cười không biết nên dùng từ ngữ gì để tả này, Kuon chỉ biết bối rối gãi đầu, mồ hôi hột tuôn ra như tắm: ``Trời đất ơi, rốt cuộc là mấy đứa đang bày ra cái trò vô tri gì ở giữa phòng họp thế này hả chứ?''

Nàng Đeo kính A-chan bấy giờ cũng rơi vào trạng thái bất lực hoàn toàn, cô nàng cố gắng vung vẩy hai cánh tay lên cao để thu hút sự chú ý của mọi người: ``Làm ơn làm phước tập trung lắng nghe tôi nói một câu đi cái.''

Thế nhưng dẫu cho cô nàng trợ lý có cố gắng ra sức lên tiếng thuyết phục, khuyên răn đến rát cả cổ họng đi chăng nữa, thì cả ba cô nàng vẫn cứ thản nhiên ngó lơ, tiếp tục tập trung toàn bộ tâm trí vào việc trêu đùa chú cừu bông của thành viên thế hệ thứ nhất, khiến cho cả A-chan lẫn Kuon đứng nhìn chỉ biết thở dài ngao ngán bất lực trước độ chai lỳ của họ.

\ct

Quyết tâm tìm ra câu chào vô địch, cả ba cô gái lại tiếp tục hăng hái thực hiện các kiểu ``chào'' độc lạ tự chế của mình. Kẻ thì hăm hở đứng lên thực hiện câu chào, kẻ thì nghiêm khắc đưa ra những lời nhận xét phê bình, rồi kẻ thì lại tỏ ra chán nản ngao ngán khi\dots\ mọi người bắt đầu có dấu hiệu cạn kiệt ý tưởng, lười suy nghĩ nên đã chuyển sang cái trò lén lút lấy luôn cái kiểu chào thương hiệu của nhau, thậm chí là ngang nhiên đi đạo nhái lấy kiểu chào của các thành viên khác trong công ty để biến tấu làm kiểu chào của riêng mình.

\begin{figure}[H]
  \centering
  \begin{subfigure}[t]{0.45\textwidth}
    \centering
    \includegraphics[width=8cm]{../../../../Image/Book4/4-13c6.png}
    \caption{Ayame bắt chước kiểu chào của Suisei\dots}
  \end{subfigure}
  \quad
  \begin{subfigure}[t]{0.45\textwidth}
    \centering
    \includegraphics[width=8cm]{../../../../Image/Book4/4-13c7.png}
    \caption{Matsuri bắt chước kiểu chào của Marine\dots}
  \end{subfigure}
  \quad
  \begin{subfigure}[t]{0.45\textwidth}
    \centering
    \includegraphics[width=8cm]{../../../../Image/Book4/4-13c8.png}
    \caption{Và Rushia lấy hình tượng trong Pok*mon\dots}
  \end{subfigure}
  \quad
  \begin{subfigure}[t]{0.45\textwidth}
    \centering
    \includegraphics[width=8cm]{../../../../Image/Book4/4-13c9.png}
    \caption{Và Ayame bắt chước kiểu chào từ Rushia.}
  \end{subfigure}
  \caption{Những kiểu chào kỳ quặc của ba thành viên, được ``mượn ý tưởng'' từ những nguồn khác nhau.}
\end{figure}

Lại đến lượt cô nàng Lễ hội bước ra sân thực hiện câu chào của mình. Cô nàng hếch cằm ngước mắt nhìn lên cao, hai ngón tay khéo léo tạo hình chữ V điệu đà đặt sát ngay bên khóe mắt nháy mắt chào hỏi: ``Chào mọi người nhé! Viên ngọc quý giá ẩn giấu trong bãi đá, ngôi sao băng rực rỡ của toàn thể mọi người đã xuất hiện rồi đây!''

Ayame đứng bên cạnh lập tức nheo mắt lên tiếng nhận xét phê bình ngay tắp lự: ``Á à! Matsuri-senpai vừa mới ngang nhiên đạo nhái trắng trợn câu chào thương hiệu của Sui-chan rồi nhé!''

Kuon đang ngồi uống trà cũng toát hết cả mồ hôi hột lo lắng giùm cô nàng: ``Anh nói thật lòng nhé Maat-chan, cái màn đạo nhái này mà lỡ chẳng may để cho Sui-chan nghe thấy được là con bé nhất định sẽ nổi trận lôi đình vác rìu sang hỏi tội em ngay lập tức đấy\dots''

Thế nhưng nàng Cổ động vẫn cười hì hì vô cùng vui vẻ, hoàn toàn chẳng thèm mảy may bận tâm lo sợ đến việc bản thân vừa mới đi đạo nhái ăn cắp ý tưởng của người khác, gương mặt vẫn toát lên một vẻ vô cùng tự hào và tự mãn.

\ct

Đến lượt trình diễn của nàng Oni. Cô nàng khẽ ngả nghiêng bả vai hướng lên bầu trời, tạo dáng vô cùng gợi cảm rồi dõng dạc cất giọng chào hỏi: ``Ahoy! Thuyền trưởng Nakiri Ayame đã có mặt rồi đây mọi người ơi!''

Matsuri lập tức lên tiếng phàn nàn chê bai: ``Cái câu chào đó rõ ràng cũng là đi đạo nhái ăn cắp ý tưởng từ cô nàng Hải tặc Marine chứ có cái gì mới mẻ đâu chứ!''

Kuon đứng bên cạnh khẽ đưa tay lên xoa xoa sống mũi của mình, mệt mỏi lắc đầu ngao ngán khuyên can: ``Anh xin hai đứa đấy, làm ơn làm phước đừng có dại dột mà réo gọi tên của Senchou lên đây nữa giùm cái đi\dots''

Nàng Quỷ sừng khẽ nở một nụ cười hờ hững đầy tinh nghịch, đôi mắt chớp chớp tỏ vẻ vô cùng ngây thơ vô tội hệt như không hề liên quan đến mình vậy.

\ct

Đến lượt của cô đồng Rushia biểu diễn. Cô nàng ngước mắt nhìn lên trần nhà, gương mặt cố gắng rặn ra một biểu cảm vô cùng giận dữ hắc ám, hai cánh tay dang rộng ra hai bên tạo dáng dọa dẫm, rồi hét lớn chào hỏi: ``Garchomp!''

Nàng Cổ động Matsuri liền vô cùng thích thú vỗ tay khen ngợi nhiệt liệt: ``Ồ! Cái ý tưởng này nghe có vẻ vô cùng độc đáo, mới mẻ và hầm hố đấy chứ. Giỏi lắm Rushia-chan!''

Thế nhưng cậu Tổng quản lý Kuon lại lập tức lên tiếng dập tắt ngay cái sự ảo tưởng đó, khẽ nhíu mày phản bác: ``Ủa khoan đã nào mấy đứa ơi, chẳng phải cái từ khóa `Garchomp' đó rõ ràng là tên của một con quái thú rồng hệ đất mà em lấy đạo nhái từ trong bộ phim hoạt hình Pok*mon ra đấy à? Sáng tạo cái nỗi gì ở đây.''

Nàng Cô đồng nghe thấy lời bóc mẽ của cậu Tổng liền lập tức xị mặt xuống, gương mặt toát lên một vẻ vô cùng bất mãn và dỗi hờn tột độ.

\ct

Nàng Quỷ sừng Ayame lại một lần nữa hào hứng lao vào giữa phòng để trổ tài. Lần này, trên tay cô nàng có cầm theo một sợi chỉ đỏ mảnh thắt nút buộc chặt lấy một đồng xu bằng đồng cổ xưa, cô nàng đứng lắc qua lắc lại sợi chỉ trước mặt mọi người, cất giọng chào hỏi: ``Hypno!''

Matsuri nheo mắt tỏ vẻ vô cùng khó hiểu và chán nản với cái cách chào hỏi lặp đi lặp lại này: ``Ủa, chẳng phải cái ý tưởng này nghe có vẻ giống hệt như cái trò chơi thôi miên tâm linh ban nãy em vừa mới phối hợp biểu diễn với Rushia-chan hay sao? Đi vòng quanh một hồi lại quay về đạo nhái chính mình à em?''

Kuon khẽ thở dài một hơi đầy bất lực trước độ vô tri của các thần tượng: ``Haiz, thế này thì cái cuộc thi vô địch Vua chào hỏi ban đầu của các em đã chính thức biến tướng trở thành cái cuộc thi `bắt chước, đạo nhái câu chào của nhau' mất rồi còn đâu\dots''

Nàng Oni nghe thấy thế liền chẳng những không hề cảm thấy xấu hổ, trái lại cô nàng còn nở một nụ cười vô cùng gian xảo đầy ranh mãnh hệt như thể bản thân vừa mới nhận được những lời khen ngợi có cánh từ phía hai người họ vậy.

\ct

Sau một khoảng thời gian dài đằng thẵng trôi qua ngập tràn trong những tiếng cười đùa rộn rã và vô vàn những câu chào tự chế vừa kỳ lạ vừa ngớ ngẩn vô tri của bộ ba Matsuri, Ayame và Rushia, toàn bộ nhóm đấu thủ rốt cuộc vẫn chẳng thể nào tìm kiếm nổi ra được một cái câu chào hỏi nào thực sự ưng ý và lọt tai cả.

Nàng Quỷ sừng chán nản nhăn mặt lại, bĩu môi tỏ vẻ vô cùng bất mãn: ``Haiz, chán quá đi mất, nghĩ nãy giờ rát cả cổ họng mà vẫn chẳng có cái câu chào nào nghe có vẻ ổn áp và dùng được cả hà\dots''

Kuon gật gù đồng cảm sâu sắc, cậu uể oải ngả ngửa người tựa hẳn ra sau thành ghế sô-pha, cất chất giọng mệt mỏi than thở: ``Thì anh đã bảo từ đầu rồi mà. Mấy cái kiểu chào tự chế của các em nhìn chung thì trông đúng là vô cùng dị hợm và độc lạ thật đấy, nhưng mà xét về tính thực tế thì đúng là nghe nó lổn nhổn, hãm tài vãi cả ra luôn ấy chứ.''

Nàng Lễ hội Matsuri vẫn quyết tâm giữ vững tinh thần chiến đấu quyết không chịu chấp nhận đầu hàng thất bại dễ dàng như vậy, cô nàng liền lập tức đưa ra một lời khuyến nghị mới: ``Thế còn cái câu chào tạm biệt kiểu từ lóng `-ò' (right-o) thì sao nào mọi người? Nghe cũng vui tai đấy!''

Rushia khẽ nhăn mặt nhíu mày phản đối ngay lập tức bằng một chất giọng vô cùng mệt mỏi, rã rời: ``Thôi đi chị ơi, nghe cái từ ngữ đó nó không được ổn định một chút nào cả-ò desu\dots''

Kuon cũng gật đầu lia lịa đồng ý với ý kiến của nàng Pháp sư: ``Anh cũng ngàn lần đồng ý-ò với ý kiến của Ru-chan luôn nhé.''

Nhận thấy cái sự không đồng tình phản đối dữ dội từ phía hai người đồng hành, nàng Lễ hội Matsuri liền tức tối nhăn mặt lại, dùng ngón tay chỉ thẳng về phía hai người họ, hắng giọng lớn tiếng ra lệnh đầy dứt khoát: ``Giời ạ, bực cả mình! Đã thế thì, lên đi, Hypno! Hãy thi triển phép thuật trừng trị tụi nó đi em!''

Chỉ chờ có thế, đôi mắt màu vàng của nàng Quỷ sừng lập tức sáng rực lên đầy hứng khởi ranh mãnh. Cô nàng nhanh nhẹn thọc tay vào túi áo lấy ra sợi chỉ buộc đồng xu đạo cụ ban nãy, bắt đầu nhịp nhàng lắc lư, đung đưa sợi chỉ qua lại liên tục ngay trước mặt mọi người. Cô nàng khẽ đanh mặt lại, nghiêm giọng lẩm nhẩm giả vờ như một pháp sư thôi miên thực thụ: ``Tất cả mọi người có mặt ở đây hãy mau chóng tập trung tinh thần, hướng mắt nhìn theo cái đồng xu của Hypno này\dots\ và từ từ chìm vào giấc ngủ sâu đi nào\dots''

Ban đầu tưởng chỉ là một trò đùa vô hại, thế nhưng một sự kiện kỳ diệu bỗng chốc xảy ra. Đôi mắt đang mở to của Rushia bỗng chốc mất dần đi ánh sáng linh hoạt thường ngày, rồi bắt đầu quay mòng mòng liên tục hệt như đang thực sự bị cái chuyển động đung đưa của chiếc đồng xu kia thôi miên hoàn toàn. Một luồng tà khí ma thuật màu tím nhạt vô cùng ảo diệu bắt đầu tỏa ra hừng hực bọc quanh cơ thể đang không ngừng đung đưa đung đưa theo nhịp của cô nàng Pháp sư.

\begin{figure}[H]
  \centering
  \begin{subfigure}[t]{0.45\textwidth}
    \centering
    \includegraphics[width=8cm]{../../../../Image/Book4/4-13d1.png}
    \caption{Matsuri chỉ tay về phía Rushia, ra lệnh cho Ayame thuyết phục cô nàng.}
  \end{subfigure}
  \quad
  \begin{subfigure}[t]{0.45\textwidth}
    \centering
    \includegraphics[width=8cm]{../../../../Image/Book4/4-13d2.png}
    \caption{Ayame lấy ra đồng xu để thôi miên Rushia.}
  \end{subfigure}
  \quad
  \begin{subfigure}[t]{0.45\textwidth}
    \centering
    \includegraphics[width=8cm]{../../../../Image/Book4/4-13d3.png}
    \caption{Rushia bị chiếc đồng xu kia điều khiển.}
  \end{subfigure}
\end{figure}

Thế nhưng điều kỳ quái và dở khóc dở cười nhất ở đây là cái ma pháp thôi miên của nàng Quỷ sừng dường như lại có sức mạnh lan tỏa quá đà. Không chỉ có mỗi một mình nạn nhân Rushia bị dính bùa, mà ngay cả thủ phạm Ayame cùng với hai người đứng xem là Matsuri và cậu Tổng Kuon đang ngồi ở bàn bấy giờ cũng cảm thấy hai mí mắt trĩu nặng lại hoàn toàn. 

Chỉ một thoáng sau đó, cả bốn con người họ đều lờ đờ, lần lượt gục đầu xuống mặt bàn họp ấm áp mà chìm sâu vào một giấc ngủ vô cùng ngon lành, sướng khoái. Cả căn phòng họp lớn bỗng chốc rơi vào một sự tĩnh lặng tuyệt đối, xung quanh giờ đây chỉ còn vang vọng đều đều tiếng thở nhẹ nhàng của bốn con người đang chìm sâu vào giấc mộng đẹp\dots

\ct

Đến khi những tia ánh sáng ban mai dịu dàng của ngày hôm sau khẽ khàng xuyên qua tấm cửa kính cửa sổ chiếu rọi vào bên trong phòng họp, cô nàng Pháp sư Rushia - người nãy giờ đang nằm sấp mặt ngủ sải lai dưới sàn nhà - mới là người đầu tiên khẽ khục khịch cựa quậy tỉnh táo mở mắt thức dậy. Cô nàng đưa hai tay lên dụi dụi đôi mắt vẫn còn hơi ngái ngủ, ngơ ngác đảo mắt nhìn quanh quất khắp phòng một lượt rồi lẩm bẩm đầy hoang mang: ``Ủa cái gì thế này\dots\ mới đó mà trời đã sáng rực rỡ từ bao giờ rồi thế này hả-nanodesu?''

\img{4-13e}

Nằm sát ngay bên cạnh cô nàng, Matsuri và Ayame bấy giờ cũng vừa mới lơ mơ bừng tỉnh giấc, cả hai cô nàng cũng đang nằm ở cái tư thế nằm sấp mặt xụi lơ dưới mặt thảm y chang như vậy, trong khi cậu Tổng quản lý Kuon thì cũng khẽ rên rỉ uể oải ngẩng đầu dậy từ trên mặt bàn làm việc chất đầy sổ sách của mình.

Cô nàng Lễ hội Matsuri đưa tay lên ngáp dài một hơi thật sảng khoái, khẽ dùng tay gãi gãi mái tóc rối bù xù xù lộn xộn của mình, ngơ ngác tự hỏi: ``Quái lạ thật sự luôn ấy chứ, tự dưng tối hôm qua bọn mình tụ tập ở đây là để định cùng nhau làm cái công việc đại sự gì thế mọi người nhỉ\dots\ sao tự dưng lại lăn ra sàn ngủ say sưa thế này không biết nữa.''

Nàng Quỷ sừng Ayame hoang mang nhìn qua nhìn lại gương mặt của mọi người một lượt, rồi lúng túng lắc lắc đầu chịu thua: ``Em hình như cũng chỉ mang máng nhớ là hình như tụi mình đang cùng nhau ngồi suy nghĩ ra ý tưởng cho một cái cuộc thi hệ trọng gì đó thôi hà\dots\ chứ cụ thể nội dung cuộc thi là cái gì thì hiện tại trong đầu em hoàn toàn trống rỗng không nhớ nổi một chữ nào nữa rồi ạ.''

Cậu Tổng quản lý Kuon khẽ đưa tay lên ôm lấy vầng trán vẫn còn hơi ê nhức mỏi mệt, nhắm nghiền đôi mắt lại như đang cố gắng vắt óc suy nghĩ lục lọi lại miền ký ức để tìm câu trả lời: ``Anh chịu chết\dots\ giờ anh cũng chả thể nào nhớ nổi ra được cái mục đích ban đầu của cái buổi họp tối hôm qua của tụi mình là để làm cái gì nữa rồi đây này\dots\ đầu óc trống rỗng hoàn toàn luôn rồi.''

Đúng vào cái khoảnh khắc cả bốn con người đang vô cùng hoang mang lú lẫn tột độ đó, cánh cửa phòng họp bỗng nhiên được đẩy nhẹ ra và cô nàng trợ lý A-chan thong thả bước chân vào phòng. Nhìn thấy cái cảnh tượng bốn con người bệ rạc vừa mới lồm cồm bò dậy từ dưới sàn nhà và mặt bàn, cô nàng trợ lý khẽ mỉm cười nhẹ nhàng gật đầu chào hỏi theo đúng thủ tục: ``Chào buổi sáng tốt lành nhé mọi người.''

Cả bốn con người họ nghe thấy tiếng chào liền bỗng chốc khựng cả người lại, đưa mắt ngơ ngác nhìn cô nàng vừa mới đặt chân vào phòng, rồi không hẹn mà gặp cùng đồng loạt cất giọng thều thào đáp lễ với một chất giọng vẫn còn đang vô cùng ngái ngủ, uể oải: ``Dạ, chúng em xin kính chào buổi sáng (chị) ạ\dots''

Có vẻ như ở thời điểm hiện tại, hoàn toàn chẳng còn sót lại lấy dù chỉ là một người trong số họ còn có khả năng ghi nhớ nổi ra được cái sự việc cái cuộc thi ngớ ngẩn mà họ đã dày công bày trò nghịch ngợm suốt cả đêm hôm trước nữa. Một cái cuộc thi vô cùng trọng đại, tưởng chừng như đã có đủ sức mạnh to lớn để thay đổi cả một cái thói quen sinh hoạt lâu đời của toàn bộ văn phòng, thế nhưng cuối cùng tất cả mọi mớ hỗn độn ấy đã chính thức bị chìm sâu vào dĩ vãng lãng quên mất rồi\dots

%==========================================TRUYỆN PHỤ=========================================
\omk

Vào buổi chiều ngày hôm đó, khi những tia ánh nắng xuân dịu dàng, ấm áp khẽ chiếu rọi qua khung cửa sổ kính của văn phòng mới, cậu con trai Kuon thong thả dọn dẹp sạch sành sanh đống sổ sách giấy tờ trên bàn họp bỏ vào cặp xách, chuẩn bị rảo bước ra về nhà sau một ngày dài làm việc căng thẳng và một đêm ngủ gục mệt mỏi ở công ty. Bước chân ra đến sát bên cạnh cánh cửa gỗ ra vào, cậu khẽ ngoảnh đầu lại nhìn ba cô gái, mỉm cười vẫy vẫy tay chào tạm biệt: ``Được rồi, công việc ngày hôm nay thế là ổn thỏa hết rồi nhé. Anh xin phép đi về nhà trước đây nha mấy đứa!''

Matsuri, Rushia và Ayame nghe thấy tiếng chào liền vô cùng vui vẻ đồng thanh cất giọng cất tiếng đáp lễ, hăm hở vẫy vẫy tay chào tạm biệt người sếp đáng kính của mình: ``Dạ vâng ạ! Matta ne(-nanodesu)! (Tạm biệt anh và hẹn gặp lại anh vào ngày mai nhé!)'' Cánh cửa gỗ phòng họp khép lại sau lưng bóng dáng của cậu Tổng quản lý, để lại ba cô nàng thần tượng tiếp tục nán lại bên trong căn phòng họp nhỏ.

Thế nhưng ngay giây tiếp theo sau khi Kuon vừa mới rời đi, nàng Quỷ sừng khẽ đưa ngón tay lên vuốt vuốt cằm ra chiều đăm chiêu suy nghĩ hệt như vừa mới phát hiện ra một điều gì đó vô cùng bất thường, cô nàng bất chợt cất tiếng hỏi: ``Cơ mà này hai chị ơi, tự dưng em lại cảm thấy cái thói quen chào hỏi tạm biệt của tụi mình có cái gì đó nghe vô cùng lạ lùng và bất hợp lý lắm luôn ấy\dots''

Nàng Lễ hội Matsuri nghiêng đầu tỏ vẻ vô cùng tò mò nhìn cô em: ``Lạ lùng ở cái chỗ nào cơ hả Ojou? Chị thấy câu chào Matta ne quen thuộc bình thường có làm sao đâu.''

Nàng Pháp sư Rushia cũng đặt nhẹ một ngón tay lên cằm ra chiều suy ngẫm trầm ngâm, cất giọng tiếp lời giải thích: ``Thì đúng là như thế thật mà, rõ ràng là tụi mình đều tự biết chắc chắn một điều rằng ngày mai bước chân vào văn phòng đi làm là bọn mình sẽ lại có cơ hội gặp mặt lại anh ấy ngay cơ mà, đúng chứ?''

``Đúng y chóc luôn đấy chị Matsuri ạ!'' Ayame gật đầu lia lịa đồng tình. ``Thế thì tại sao hàng ngày tụi mình vẫn cứ phải giữ cái thói quen dùng cái câu chào tạm biệt Matta ne nghe nó cứ có vẻ xa cách, khách sáo hệt như là cả đời này tụi mình không có cơ hội gặp lại nhau nữa ấy\dots\ nghe nó vô lý đùng đoàng ra luôn ấy chứ, yo.''

Matsuri nhún vai một cái nhẹ tênh, khẽ cười nhẹ đáp lại: ``Ối dào cái con bé này vẽ chuyện quá cơ, chị thấy cái thói quen đó là hoàn toàn bình thường mà. Ở cái đất nước Nhật Bản này xưa nay mọi người ai ai cũng đều dùng cái câu chào tạm biệt cửa miệng quen thuộc đó để tiễn nhau ra về thôi chứ có ai thắc mắc gì đâu.''

Rushia gật gù gật gù đầu như vừa mới sực nhớ ra điều gì: ``À đúng rồi nè, em nhớ mang máng là bên bộ phận phòng nhân sự người ta hồi trước đi làm cũng từng có lần căn dặn bảo với bọn mình như vậy rồi mà-nanodesu\dots''

Cô nàng trợ lý A-chan nãy giờ vẫn đang ngồi cắm cúi làm việc giấy tờ ở chiếc bàn họp gần đó, nghe thấy ba cô em bắt đầu dở chứng bàn luận xàm xí liền không nhịn được mà bất ngờ xen ngang cất giọng dứt khoát hòng chấm dứt chủ đề: ``Thì tại vì cái thói quen quy tắc chào hỏi tạm biệt từ xưa đến nay nó đã là bắt buộc như vậy rồi mà! Mấy đứa đừng có mà rảnh rỗi đứng đó vẽ chuyện bao đồng nữa đi nha!''

Thế nhưng nàng Quỷ sừng vẫn quyết tâm không chịu từ bỏ ý định quậy phá của mình. Cô nàng tinh nghịch giả vờ như điếc, hoàn toàn ngó lơ phản đối cộc lốc của đàn chị Đeo kính, khẽ nghiêng nghiêng đầu hướng mắt nhìn ra ngoài cửa sổ xa xăm: ``Hmm\dots\ Nhưng mà tại sao cái thói quen chào hỏi tạm biệt đó nó lại được hình thành như thế nhỉ? Có ai giải thích giùm em một câu được không ạ?''

Cô nàng trợ lý đeo kính A-chan khẽ thở dài bất lực, cất giọng nhắc nhở lại một lần nữa, lần này cô nàng cố tình nhấn mạnh rõ ràng từng từ ngữ hệt như đang dạy trẻ con mẫu giáo học chữ: ``Thì tôi đã bảo là cái thói quen quy tắc chào hỏi đó nó bắt buộc phải như vậy rồi mà! Cấm có cãi bướng nha cái cô kia!''

Thế nhưng, ba cô nàng thần tượng nghịch ngợm bấy giờ bỗng chốc im lặng trầm ngâm mất vài giây, mỗi người tự chìm đắm vào một dòng suy nghĩ tưởng tượng riêng biệt của bản thân mình. Chỉ vài giây ngắn ngủi sau đó, cô nàng Lễ hội Matsuri bỗng chốc reo hò lên đầy phấn khích như thể vừa mới có một tia chớp ý tưởng sáng tạo tuyệt vời rạch ngang qua đại não: ``A! Chị nghĩ ra ý kiến này hay ho lắm nè hai đứa ơi! Tại sao ba chị em tụi mình không thử cùng nhau động não, sáng tạo ra một cái cách thức chào hỏi tạm biệt mới mẻ độc lạ mang đậm nét đặc trưng riêng biệt của văn phòng tụi mình đi xem nào?''

Ojou-san Ayame nghe thấy thế liền lập tức gật đầu cái rụp đồng ý ngay tắp lự, đôi mắt sáng rực lên rạng rỡ hệt như vừa mới tìm thấy được một rương kho báu vô giá: ``Ý kiến đề xuất này nghe có vẻ vô cùng thú vị và tuyệt vời ông mặt trời lắm đó! Em duyệt cái phương án này ngay lập tức luôn nhé!''

Trái ngược hoàn toàn với cái sự phấn khích tột độ của hai người đồng nghiệp, nàng Pháp sư Rushia lại khẽ nhíu mày tỏ vẻ vô cùng e sợ và lo ngại sâu sắc: ``Nhưng mà\dots\ mọi người làm ơn hãy tự nhìn lại cái tấm gương tàn tạ của ngày hôm qua đi ạ, chẳng phải cũng chính vì mấy cái ý tưởng sáng tạo câu chào dở hơi kiểu này mà tụi mình đã phải trả giá đắt rồi hay sao-nanodesu\dots''

Cô nàng Đeo kính A-chan ngồi bên cạnh đưa mắt chứng kiến cái cảnh tượng ba cô em gái lại bắt đầu hớn hở chụm đầu vào nhau để vẽ chuyện suy nghĩ ra những câu chào tạm biệt đầy kỳ dị và ngớ ngẩn vô tri mới, chỉ biết lặng lẽ lắc đầu ngao ngán bất lực tột cùng, khẽ lẩm bẩm thở dài một câu: ``\hl{Thôi xong, cái kịch bản dở hơi này lại chuẩn bị tái diễn thêm một lần nữa rồi\dots}''

Và thế là, hoàn toàn đúng y như những gì cô nàng lo sợ, một vòng lặp của một ``cuộc thi quyết đấu'' ngớ ngẩn vô tri khác lại chính thức được khởi động bắt đầu, y hệt như cái cách mà ba cô nàng đã từng bày trò quậy phá vào đêm hôm trước; một buổi tối hứa hẹn sẽ lại tiếp tục tràn ngập sự hỗn loạn phi lý, tiếng cười đùa sảng khoái và vô vàn những cái ý tưởng chào hỏi tạm biệt không giống ai xuất hiện.

Có thể ở thời điểm hiện tại, bản thân ba cô nàng thần tượng hoàn toàn không còn giữ lại được lấy dù chỉ là một mẩu ký ức nhỏ nhặt nào về những sự kiện tai hại, dở khóc dở cười mà họ đã từng làm vào tối ngày hôm qua, thế nhưng sâu thẳm trong tiềm thức vô tri của ba người họ bấy giờ, dường như ai nấy đều đã chuẩn bị sẵn sàng tâm lý rực lửa để chào đón những câu chào tạm biệt vô cùng ``trang trọng'' và ``tuyệt hảo'' nhất cho cuộc thi sắp sửa nổ ra phía trước.

\begin{center}{\abv\Large BONUS GIF}\end{center}

\gif{4-13f}{1}{29}

\end{document}